% !TeX root = ../Book2.tex
\chapter*{后记}
\addcontentsline{toc}{chapter}{后记}
读到本卷末尾时，一个空间已经可以通过几种不同的资料来研究：胞腔与附着映射、链复形与上同调环、映射的同伦类、临界点和把手，以及其上的向量丛。每一种资料都使某些问题更容易计算。把它们联系起来的比较定理，说明何时可以换一种描述继续推理，也说明转换时保留了哪些信息。

\ProseParagraphBreak

映射度贯穿了这样的转换。它在胞腔边界中记录附着的次数，在光滑映射中成为逆像点的带符号计数，又通过积分与微分形式相连。Euler 类也把几种原本相距较远的构造放到一起：向量丛的零截面、横截截面的零点、上同调中的类，以及合适几何条件下的曲率表达。回头把这些构造中的定向和自然性逐一对齐，可以更清楚地理解同一个不变量为何有多种计算方法。

\ProseParagraphBreak

本卷为链级构造、几何实现和比较证明保留了较多篇幅。连接映射如何选取代表，交点的符号如何确定，局部系数如何随路径移动，都会影响最后的答案。第一次阅读时，这些细节可能使证明显得缓慢；真正计算一个带扭系数的例子，或比较两个文献中的定向约定时，它们便成为能够直接使用的依据。重读时可以挑选一次完整计算，追踪每项选择进入公式的位置，并检查更换选择后发生的变化。

\ProseParagraphBreak

滤过和稳定化也带来了新的阅读方式。谱序列将一个复杂计算拆成逐页修正，最后还要从滤过的相邻商恢复整个群及其乘法结构；稳定方法允许添加平凡方向，从而形成更便于运算的对象。使用这些方法时，适合把所求对象和实际算出的资料并排写清楚。比如，一个上同调环、一组示性数和一个稳定丛类所回答的问题各不相同；进一步分类时，需要知道还有什么几何信息尚未被记录。

\ProseParagraphBreak

后续几卷将继续使用本卷的工具。临界点与变分问题相接，横截性进入几何对象的参数空间，示性类和曲率代表出现在复几何与规范几何中，K 理论中的符号类则连接拓扑与椭圆分析。第三卷将展开与几何方程有关的分析方法，第四卷研究复、辛结构及其相容条件，第五卷在代数几何中发展层上同调、相交与模问题。相近的名称往往伴随新的对象和条件，学习时可以以本卷的构造作参照，再核对新的定义与比较结果。

\ProseParagraphBreak

进一步阅读可以从自己仍想算清的例子开始。若关心纤维化，可以继续追踪谱序列中的微分与隐藏延拓；若关心流形，可以比较不同临界点分解及其把手变化；若关心向量丛，可以先计算具体丛的示性类，再研究稳定化后的分类。低维拓扑、分层空间和有理同伦等附录提供了另外一些选择。能够把一个问题化成明确的构造，选出适用的不变量，并判断计算结果足以回答哪一部分，便有了继续学习这些专题的具体起点。
